\section{反常积分}
\label{sec:ru-improper-integrals}

黎曼积分只对有限区间上的有界函数定义。下面两个积分都不在这个范围内：$\int_1^\infty x^{-2}\,\mathrm dx$的积分区间无界；$\int_0^1x^{-1/2}\,\mathrm dx$的区间有限，被积函数却在$0$附近无界。
处理办法是先在一个有限、有界的部分上积分，再让这个部分扩大，看积分是否有极限。

\subsection{无穷区间与奇点处的截断极限}
\label{sec:ru-improper-definition}

设$f$在每个$[a,b]$（$b>a$）上可积。若$\lim_{b\to+\infty}\int_a^bf$存在且有限，就称\term{反常积分}（improper integral）$\int_a^{+\infty}f$收敛，并以这个极限为它的值；否则称发散。
例如$\int_1^bx^{-2}\,\mathrm dx=1-1/b\to1$，所以$\int_1^\infty x^{-2}\,\mathrm dx=1$；而$\int_1^bx^{-1}\,\mathrm dx=\log b\to+\infty$，$\int_1^\infty x^{-1}\,\mathrm dx$发散。
被积函数在端点附近无界时，先避开这个点再取单侧极限：若$f$在每个$[a+\varepsilon,b]$上可积，定义$\int_a^bf=\lim_{\varepsilon\to0^+}\int_{a+\varepsilon}^bf$。
例如$\int_\varepsilon^1x^{-1/2}\,\mathrm dx=2-2\sqrt\varepsilon\to2$。函数在端点附近无界的点称为奇点。

区间两端都无界，或者区间内部有奇点时，要在中间取一点把积分拆开，要求每一段分别收敛，再相加；结果不依赖拆分点的选择。
这一要求防止用抵消掩盖某一侧的发散。$\int_{-\infty}^{+\infty}x\,\mathrm dx$的对称截断$\int_{-b}^bx\,\mathrm dx$恒为零，但$\int_0^{+\infty}x\,\mathrm dx$发散，所以这个反常积分发散。
同样，$\int_{-1}^1x^{-1}\,\mathrm dx$在内部有奇点，两侧$\int_0^1x^{-1}\,\mathrm dx$与$\int_{-1}^0x^{-1}\,\mathrm dx$都发散，它也发散。
对称截断的极限称为Cauchy主值，上面两个积分的主值都是$0$。主值是另一种规定，在特定问题中有用，但不能当作通常反常积分的值。

\subsection{收敛判别、绝对收敛与尾部误差}
\label{sec:ru-improper-tests}

若$f\geq0$，$\int_a^bf$随$b$单调不减，所以$\int_a^\infty f$收敛当且仅当这些截断积分有上界，与非负项级数的情形相同。由此得到比较判别：
若$0\leq f\leq g$且$\int_a^\infty g$收敛，则$\int_a^\infty f$收敛。不要求非负时，用柯西准则：$\int_a^\infty f$收敛，当且仅当对每个$\varepsilon>0$，存在$B$使$b_2>b_1>B$时$\bigl|\int_{b_1}^{b_2}f\bigr|<\varepsilon$。
比较需要基准。由原函数直接计算，
\[
 \int_1^\infty x^{-p}\,\mathrm dx\ \text{收敛}\iff p>1,\qquad\int_0^1x^{-p}\,\mathrm dx\ \text{收敛}\iff p<1.
\]
两者方向相反：无穷远处需要衰减得足够快，零点附近则不能增长得太快。

收敛的反常积分可以用有限区间上的积分近似，误差就是剩下的尾部，比较判别同时给出它的上界。例如$\int_0^\infty e^{-x^2}\,\mathrm dx$：在$x\geq b>0$上$e^{-x^2}\leq(x/b)e^{-x^2}$，所以
\[
 0\leq\int_b^\infty e^{-x^2}\,\mathrm dx\leq\frac1b\int_b^\infty xe^{-x^2}\,\mathrm dx=\frac{e^{-b^2}}{2b}.
\]
截断到$[0,3]$时，尾部误差不超过$e^{-9}/6\approx2\times10^{-5}$。剩下的有限区间积分再用第~\ref{sec:ru-numerical-quadrature}节的求积公式计算。

\paragraph{绝对收敛与条件收敛}
若$\int_a^\infty|f|$收敛，称$\int_a^\infty f$绝对收敛。由柯西准则和$\bigl|\int_{b_1}^{b_2}f\bigr|\leq\int_{b_1}^{b_2}|f|$，绝对收敛蕴含收敛。反之不成立，$\int_1^\infty\frac{\sin x}x\,\mathrm dx$是代表例子。
对截断积分分部积分，取$u=1/x$、$v'=\sin x$，
\[
 \int_1^b\frac{\sin x}x\,\mathrm dx=\Bigl[-\frac{\cos x}x\Bigr]_1^b-\int_1^b\frac{\cos x}{x^2}\,\mathrm dx.
\]
$b\to\infty$时边界项趋于$\cos1$，右边的积分由$|\cos x|/x^2\leq1/x^2$绝对收敛，所以原积分收敛。但在每个$[k\pi,(k+1)\pi]$上$\int|\sin x|/x\geq\frac2{(k+1)\pi}$，累加起来是调和级数，所以$\int_1^\infty|\sin x|/x$发散。
这个积分的收敛来自正负部分的振荡抵消，称为条件收敛，与交错级数的情形相似。反常积分收敛与一般测度意义下的可积是否相同，要等第~\ref{chap:measurable-functions}章定义了相应的积分以后才能比较。
